
	
	
	\section{有电介质时的高斯定理和环路定理  电位移}
	
	
	\begin{definition}[电位移]
		电位移用$D$表示，单位是$C/m^2$。
		$$D = \varepsilon_{0}E + P$$
		其中，$P$为电极化强度。
	\end{definition}
	
\begin{definition}[电位移通量]
在有电介质的静电场中作电位移线，使线上每一点的切线方向和该点电位移 D 的方向相同，并规定在垂直于电位移线的单位面积上通过的电位移线数目等于该点的电位移 D 的量值，被称为称作电位移通量。
\end{definition}
	
	\begin{definition}[有电介质时的高斯定理]
		$$\oint_{S} D \cdot dS = q_{0}$$
	\end{definition}
	
\begin{note}
有电介质时的高斯定理是静电场基本定理之一，是普遍适用的。
	
通过电介质中任一闭合曲面的电位移通量，等于该面所包围的自由电荷量的代数和。

电位移线是从正的自由电荷出发，终止于负的自由电荷（不包括极化电荷）。

\end{note}
	
	
\begin{note}
$D$和$E$和$P$的关系和区分：

有电介质时，平行板电容器之间$D$线均匀分布（从正的自由电荷出发，终止于负的自由电荷）（不包括极化电荷），$E$线有部分终止于极化电荷，在电介质内较为疏松，$P$线起始于负极化电荷，终止于正极化电荷，只在电介质内部。

设电极化率为$\chi_{e}$，相对电容率$\varepsilon_{0}$，电容率$\varepsilon$
$$D = \varepsilon_{0}E + P = \varepsilon_{0}E + \chi_{e}\varepsilon_{0}E = \varepsilon_{r}\varepsilon_{0}E$$
	$$D = \varepsilon E$$
		
\end{note}
	
	\begin{definition}[有电介质时的环路定理]
$$\oint_{L} E \cdot dl = 0$$
$E$为所有电荷激发的静电场。
	\end{definition}
\begin{supplement}
带电金属球周围充满均匀无限大电介质后，其电场强度减弱到真空时的$\frac{1}{\varepsilon_{r}}$
\end{supplement}

\begin{supplement}
$E = \frac{E_{0} }{\varepsilon_{r}}$ 和 $D = \varepsilon_{0}E $仅在均匀电介质充满整个电场或电介质表面是等势面时成立。
\end{supplement}
	

	
